\begin{table}[htbp]
\centering
\footnotesize
\caption{Le destinazioni trattate: quando l'accordo entra in vigore e quanto contenuto ambientale contiene}
\label{tab:trattamento}
\begin{threeparttable}
\begin{tabular}{llcc}
\toprule
Destinazione & Codice & Anno di entrata & Profondit\`a EP massima \\
 & & in vigore & WB \quad / \quad TREND \\
\midrule
Bangladesh & 103 & 2002 & 1 \quad / \quad 1 \\
India & 111 & 2002 & 1 \quad / \quad 1 \\
Korea Rep. & 133 & 2002 & 17 \quad / \quad 72 \\
Laos,PDR & 119 & 2002 & 6 \quad / \quad 4 \\
Sri Lanka & 134 & 2002 & 1 \quad / \quad 1 \\
HongKong & 110 & 2003 & 4 \quad / \quad 5 \\
Macau & 121 & 2003 & 5 \quad / \quad 5 \\
Brunei & 105 & 2005 & 6 \quad / \quad 4 \\
Cambodia & 107 & 2005 & 6 \quad / \quad 4 \\
East Timor & 144 & 2005 & 6 \quad / \quad 4 \\
Indonesia & 112 & 2005 & 6 \quad / \quad 4 \\
Malaysia & 122 & 2005 & 6 \quad / \quad 4 \\
Myanmar & 106 & 2005 & 6 \quad / \quad 4 \\
Philippines & 129 & 2005 & 6 \quad / \quad 4 \\
Singapore & 132 & 2005 & 7 \quad / \quad 15 \\
Thailand & 136 & 2005 & 6 \quad / \quad 4 \\
Vietnam & 141 & 2005 & 6 \quad / \quad 4 \\
Chile & 412 & 2006 & 5 \quad / \quad 30 \\
Pakistan & 127 & 2007 & 3 \quad / \quad 12 \\
New Zealand & 609 & 2008 & 7 \quad / \quad 53 \\
Peru & 434 & 2010 & 12 \quad / \quad 47 \\
Costa Rica & 415 & 2011 & 4 \quad / \quad 50 \\
Iceland & 322 & 2014 & 6 \quad / \quad 28 \\
Switzerland & 331 & 2014 & 14 \quad / \quad 53 \\
Australia & 601 & 2015 & 3 \quad / \quad 24 \\
\midrule
\textbf{Totale} & 25 destinazioni & & \\
\bottomrule
\end{tabular}
\begin{tablenotes}[flushleft]\footnotesize
\item Ogni riga \`e una destinazione con cui la Cina ha un accordo commerciale in vigore fra il 2000 e il 2015.
\item \textit{Profondit\`a EP} = quante disposizioni ambientali contiene l'accordo, secondo due codifiche indipendenti: quella della Banca Mondiale (WB) e quella del database accademico TREND.
\item Le 11 destinazioni ASEAN condividono lo stesso accordo (2005) e quindi gli stessi valori: gli accordi realmente distinti sono circa 14, non 25. \`E questo il vincolo che governa tutta l'inferenza del lavoro.
\item Hong Kong e Macao sono economie di transito: met\`a del valore esportato verso destinazioni trattate passa da l\`i. Per questo il campione principale le esclude e una variante le reinserisce.
\end{tablenotes}
\end{threeparttable}
\end{table}
